\subsection{Data Sources and Experimental Roles}

UAV-SEAD supplied the supervised training, runtime detection, and fault-domain diagnosis data. ALFA was used only for zero-shot binary anomaly detection, and Uncategorized UAV-SEAD logs were used only for flight-level weak-label screening. Neither analysis expanded the four-class supervised labels. Controlled PX4 replay was kept separate from real-flight model fitting and provided the only module-level mutation ground truth. The roles, protocol sizes, model configurations, and evaluation outputs are summarized in Table~\ref{tab:protocol}.

Two recent multivariate time-series methods were also examined as detector references. CATCH \cite{wu2025catch} used the official model core at repository commit \texttt{\seqsplit{3647c69be5eb56649b072596cf89098e689e20c3}}, all 87 raw telemetry channels, and the same fixed chronological train/validation/test windows as HS-AeroTS. Following its unsupervised task definition, each of three seeds trained only on a deterministic sample of 4096 normal training windows; model selection used 1024 normal validation windows. The declared compute budget was two epochs with batch size 8; evaluation batch size 128 changed only inference throughput. This matched-input execution removes the earlier 18-topic aggregation but remains a budget-limited, normal-only comparison rather than an equal-supervision contest. A GCAD-inspired reference \cite{liu2025gcad} used a ridge-regularized linear VAR(1) Granger-prediction score on 8192 normal windows; it is not a reproduction of the published nonlinear GCAD network. These baselines are reported with their actual protocols and are not used to claim unrestricted state-of-the-art superiority.
